\documentclass[11pt, letterpaper]{article}

% Packages:
\usepackage[
    ignoreheadfoot,
    top=1.4 cm,
    bottom=1.4 cm,
    left=1.4 cm,
    right=1.4 cm,
    footskip=0.6 cm
]{geometry}
\usepackage{titlesec}
\usepackage{tabularx}
\usepackage{array}
\usepackage[dvipsnames]{xcolor}
\definecolor{primaryColor}{RGB}{0, 0, 0}
\usepackage{enumitem}
\usepackage{fontawesome5}
\usepackage{amsmath}
\usepackage[
    pdftitle={Tianjian Feng's CV},
    pdfauthor={Tianjian Feng},
    pdfcreator={LaTeX with RenderCV},
    colorlinks=true,
    linkcolor=black,
    urlcolor=primaryColor
]{hyperref}
\usepackage[pscoord]{eso-pic}
\usepackage{calc}
\usepackage{bookmark}
\usepackage{lastpage}
\usepackage{changepage}
\usepackage{paracol}
\usepackage{ifthen}
\usepackage{needspace}
\usepackage{iftex}
\usepackage{xcolor}

% Ensure that generated PDF is machine readable/ATS parsable:
\ifPDFTeX
    \input{glyphtounicode}
    \pdfgentounicode=1
    \usepackage[T1]{fontenc}
    \usepackage[utf8]{inputenc}
    \usepackage{lmodern}
\fi

\usepackage{charter}

% Some settings:
\raggedright
\AtBeginEnvironment{adjustwidth}{\partopsep0pt}
\pagestyle{empty}
\setcounter{secnumdepth}{0}
\setlength{\parindent}{0pt}
\setlength{\topskip}{0pt}
\pagenumbering{gobble}

\titleformat{\section}
{\needspace{4\baselineskip}\bfseries\normalsize}
{}{0pt}{}
[\vspace{0pt}\titlerule]

\titlespacing{\section}{
    -2pt
}{
    0.05 cm
}{
    0.05 cm
}

\renewcommand\labelitemi{\textbullet}

\newenvironment{highlights}{
    \begin{itemize}[
        topsep=0 cm,
        parsep=0 cm,
        partopsep=0pt,
        itemsep=0pt,
        leftmargin=0 cm + 10pt
    ]
}{
    \end{itemize}
}

\newenvironment{highlightsforbulletentries}{
    \begin{itemize}[
        label=\vcenter{\hbox{\textbullet}},
        topsep=0cm,
        parsep=0cm,
        partopsep=0pt,
        itemsep=0pt,
        leftmargin=0cm + 10pt,
        font=\normalsize
    ]
}{
    \end{itemize}
}

\newenvironment{onecolentry}{
    \begin{adjustwidth}{
        0 cm + 0.00001 cm
    }{
        0 cm + 0.00001 cm
    }
}{
    \end{adjustwidth}
}

\newenvironment{twocolentry}[2][]{
    \onecolentry
    \def\secondColumn{#2}
    \setcolumnwidth{\fill, 5 cm}
    \begin{paracol}{2}
}{
    \switchcolumn \raggedleft \secondColumn
    \end{paracol}
    \endonecolentry
}

\newenvironment{threecolentry}[3][]{
    \onecolentry
    \def\thirdColumn{#3}
    \setcolumnwidth{, \fill, 4.5 cm}
    \begin{paracol}{3}
    {\raggedright #2} \switchcolumn
}{
    \switchcolumn \raggedleft \thirdColumn
    \end{paracol}
    \endonecolentry
}

\newenvironment{header}{
    \setlength{\topsep}{0pt}
    \par\kern\topsep
    \centering
    \linespread{0.95}
}{
    \par\kern\topsep
}

\let\hrefWithoutArrow\href

\usepackage{setspace}
\setstretch{1.15}

\begin{document}

    \newcommand{\AND}{\unskip
        \cleaders\copy\ANDbox\hskip\wd\ANDbox
        \ignorespaces
    }
    \newsavebox\ANDbox
    \sbox\ANDbox{$|$}

    \begin{header}
        \fontsize{17 pt}{18 pt}\selectfont \textbf{Tianjian Feng}

        \normalsize
        \kern 3.0 pt%
        \mbox{\hrefWithoutArrow{mailto:tianjian@illinois.edu}{tianjian@illinois.edu}}%
        \kern 3.0 pt%
        \AND%
        \kern 3.0 pt%
        \mbox{\hrefWithoutArrow{tel:+12172550587}{+1 (217) 255-0587}}%

        \vspace{2 pt}

        \textit{J-1 Student Intern / Visiting Undergraduate Researcher,
        University of Illinois Urbana-Champaign}
    \end{header}


    \section{Education}

    \begin{twocolentry}{
        \makebox[5cm][r]{\textit{May 2026 -- Present}}
    }
        \textbf{University of Illinois Urbana-Champaign}, USA
    \end{twocolentry}

    \begin{onecolentry}
        J-1 Student Intern / Visiting Undergraduate Researcher,
        Siebel School of Computing and Data Science
    \end{onecolentry}

    \vspace{0.15 cm}

    \begin{twocolentry}{
        \makebox[5cm][r]{\textit{Sep. 2023 -- June 2027 (expected)}}
    }
        \textbf{Zhejiang University}, China
    \end{twocolentry}

    \begin{onecolentry}
        Bachelor of Computer Science;
        Cumulative GPA: 4.09/4.3 (Top 25\%)
    \end{onecolentry}

    \vspace{0.15 cm}


    \section{Research Interests}

    \vspace{0.1cm}

    My research focuses on scaling and optimizing generative foundation models across modalities.
    I am particularly interested in the intersection of structured reasoning and generative modeling,
    including efficient decoding, test-time scaling, physically consistent video generation,
    and latent representation learning.
    My goal is to develop robust, high-fidelity generative frameworks that maintain
    spatial-temporal consistency and logical integrity in complex inference tasks.


    \section{Publications and Pre-prints}

    \begin{onecolentry}
    \begin{highlights}

    \item
    \textit{(EMNLP 2026)}
    Y. Shen$^*$, \textbf{T. Feng}$^*$, J. Han, W. Wang, T. Chen,
    C. Shen, J. Leskovec, and S. Ermon. (2026).
    Improving Diffusion Language Model Decoding through Joint Search in
    Generation Order and Token Space.
    \href{https://arxiv.org/abs/2601.20339}
    {\textcolor{NavyBlue}{[paper]}}
    \textit{* Equal contribution.}

    \item
    \label{pub:Tinker}
    \textit{(ICLR 2026)}
    C. Zhao, X. Li, \textbf{T. Feng}, Z. Zhao, H. Chen, and C. Shen. (2026).
    TINKER: Diffusion's Gift to 3D--Multi-View Consistent Editing From Sparse
    Inputs without Per-Scene Optimization.
    \href{https://arxiv.org/abs/2508.14811}
    {\textcolor{NavyBlue}{[paper]}}
    \href{https://aim-uofa.github.io/Tinker/}
    {\textcolor{NavyBlue}{[project]}}

    \item
    \textit{(ACL 2026)}
    L. Zhong$^*$, L. Wu$^*$, B. Fang, \textbf{T. Feng}, C. Jing,
    W. Wang, J. Zhang, H. Chen, and C. Shen. (2026).
    Beyond Hard Masks: Progressive Token Evolution for Diffusion Language Models.
    \href{https://arxiv.org/abs/2601.07351}
    {\textcolor{NavyBlue}{[paper]}}
    \href{https://aim-uofa.github.io/EvoTokenDLM/}
    {\textcolor{NavyBlue}{[project]}}

    \item
    \textit{(ECCV 2026)}
    L. Li, W. Wang, C. Zhao, \textbf{T. Feng}, Z. Zhao, H. Chen, and C. Shen.
    (2026).
    MMControl: Unified Multi-Modal Control for Joint Audio-Video Generation.
    \href{https://arxiv.org/abs/2604.19679v2}
    {\textcolor{NavyBlue}{[paper]}}
    \href{https://aim-uofa.github.io/MMControl/}
    {\textcolor{NavyBlue}{[project]}}
    \href{https://github.com/aim-uofa/MMControl}
    {\textcolor{NavyBlue}{[code]}}

    \end{highlights}
    \end{onecolentry}


    \section{Research Experiences}

        \begin{onecolentry}
            \textbf{University of Illinois Urbana-Champaign} \hfill May 2026 -- Present \\
            \textit{Research Intern (Advisor: Prof. Yuxiong Wang)}

            \begin{highlights}

                \item \textbf{Physics-Aware Video Continuation} -- Ongoing
                \begin{itemize}
                    \item Developing a prefix-conditioned video continuation pipeline
                    to study how structured predictions of scene dynamics can guide
                    future-frame generation.

                    \item Implementing teacher-supervised fine-tuning of a
                    vision-language reasoner using full-window video annotations,
                    while restricting its inputs to observed prefixes and
                    prefix-derived captions.

                    \item Conducting controlled experiments to distinguish errors in
                    prefix perception, future-event prediction, and video synthesis,
                    using matched generation settings and event-level evaluation.
                \end{itemize}

            \end{highlights}
        \end{onecolentry}

        \vspace{0.15 cm}

        \begin{onecolentry}
            \textbf{State Key Lab of CAD \& CG, Zhejiang University}
            \hfill 2024 -- Present \\

            \textit{Undergraduate Researcher
            (Advisor: Prof. Chunhua Shen, Research Prof. Hao Chen)}

            \begin{highlights}

                \item \textbf{Efficient Decoding for Diffusion Language Models (DLM)}
                \begin{itemize}
                    \item \textbf{Spearheaded the technical implementation}
                    of the Order-Token Search algorithm and its core likelihood estimator,
                    ensuring efficient pruning and stable exploration across diverse
                    decoding trajectories.
                \end{itemize}

                \item \textbf{High-Fidelity 3D Content Generation and Editing}
                \begin{itemize}
                    \item Contributed to \textbf{TINKER}, a generalizable 3D editing
                    framework that achieves state-of-the-art performance without
                    per-scene optimization \textit{(ICLR 2026)}.

                    \item Conducted \textbf{extensive benchmarking} to evaluate the
                    framework's performance in one-shot and few-shot regimes,
                    providing critical experimental evidence for model validation.
                \end{itemize}

            \end{highlights}
        \end{onecolentry}


    \section{Skills}

        \begin{onecolentry}
            \textbf{Programming:} Python, C++, Java, Verilog, Shell, \LaTeX; \\
            \textbf{Frameworks:} PyTorch, Verl.
        \end{onecolentry}


    \section{Honors \& Awards}

        \begin{onecolentry}
            \textbf{Zhejiang University Academic Excellence Scholarship}
            \hfill 2024, 2025 \\

            \textit{Awarded to the top 20\% of students for outstanding
            academic and research performance.}
        \end{onecolentry}

\end{document}
